\chapter{Occupation, Not Loudness}
\label{ch:intro}

\section{A familiar complaint, stated precisely}

Few complaints about home viewing are as common as the one about dialogue and explosions. The volume that makes speech intelligible makes the action unbearable; the volume that makes the action tolerable makes the speech inaudible. Viewers respond by keeping a hand on the remote. Industry responses have been largely framed in terms of loudness: broadcast loudness normalisation, legislation against loud advertisements, and dynamic range compression modes built into consumer decoders. These responses treat the problem as one of level.

This monograph takes a different view. The difficulty is not that a soundtrack is loud, nor even that it has a wide dynamic range. Both are ordinary properties of cinema sound, and both are intended. The difficulty arises from three conditions holding together:
\begin{enumerate}
\item the listening level is set for dialogue, because dialogue is what must be understood;
\item non-dialogue material exceeds that level by a large margin; and
\item it does so for long, nearly continuous stretches, with few intervals in which the level returns to the dialogue range.
\end{enumerate}
The first condition fixes a reference. The second is a question of magnitude. The third is a question of temporal organisation, and it is the one that loudness measures summarise away. An integrated loudness figure for a film is a single number over its whole running time; a loudness range figure describes a distribution. Neither says whether the loud material arrives as occasional peaks separated by long rests or as a continuous block several minutes long. The distinction between those two cases is the subject of this book.

\section{Peak intensity and occupation}

The central distinction can be stated compactly. \emph{Peak intensity} concerns how far above the reference a soundtrack goes. \emph{Occupation} concerns how much of the time it stays there, and how continuously. A conventional dramatic structure has a rhythm of rest, rise, peak, and release. A soundtrack can instead proceed through rise, impact, barrage, further rise, and further impact without returning to rest. The second pattern need not have higher peaks than the first. What distinguishes it is that the listener is never released.

The practical consequence is that the listener's volume control becomes part of the soundtrack's reproduction. If the reference is set for speech and the soundtrack occupies the range well above it for minutes at a time, the listener must either accept the exceedance, reduce the volume and lose the next line of dialogue, or reduce the volume to nothing and wait. The case at the centre of this study involved the last of these. For action sequences lasting one to two minutes or more, the volume was turned fully down until the scene calmed.

\section{Why a single case}

This monograph is built around one film, \emph{Drive Through Fire}, and one reproduction setup. That is a limitation and it is stated as such. It is also a methodological choice. The concepts needed to describe the problem (dialogue reference, exceedance, occupation block, recovery window, salience channel, throttle) are not yet standard, and the instruments needed to measure them do not exist off the shelf. A single case examined closely is the right scale at which to build and debug those instruments. The case is not offered as evidence that contemporary film mixing is generally at fault. It is offered as a worked example of how that question could be answered.

\section{Grading claims}

Every substantive claim in the book is assigned one of four grades, and Chapter~\ref{ch:conclusion} tabulates them.

\begin{description}
\item[Established.] Claims supported by standards, published measurement, or well-replicated psychoacoustics, independent of this study. Example: cinema reproduction is calibrated to a reference sound pressure level; home reproduction generally is not.
\item[Measured.] Claims about the \emph{Drive Through Fire} soundtrack that follow from the analysis pipeline as corrected, with the caveats of Chapters~\ref{ch:pipeline} and~\ref{ch:failures}.
\item[Hypothesised.] Claims that the framework predicts and that the measurements are consistent with, but that require listener data, a comparison corpus, or measurement of the playback chain to test.
\item[Rejected.] Claims that were made during the analysis, sometimes confidently, and that did not survive inspection. These are kept because the errors are instructive.
\end{description}

\section{Plan of the book}

Chapter~\ref{ch:background} surveys the standards and literatures that bear on the problem: loudness measurement, cinema calibration, film sound practice and theory, psychoacoustics, and noise annoyance. Chapter~\ref{ch:framework} gives formal definitions and states the hypotheses. Chapter~\ref{ch:pipeline} describes the analysis pipeline. Chapter~\ref{ch:failures} documents its failures and their corrections. Chapter~\ref{ch:chain} treats the reproduction chain as an object of study in its own right. Chapter~\ref{ch:case} presents the case. Chapter~\ref{ch:overdetermination} examines the co-occurrence of salience channels. Chapter~\ref{ch:response} specifies the response measure and the study design needed to test the framework against listener behaviour. Chapter~\ref{ch:intervention} develops an intervention. Chapter~\ref{ch:conclusion} grades the results.
